\section{Conclusions}

This study implemented and compared a plain Gaussian EDA, an EDA with tightly budgeted L-BFGS-B refinement, and an EDA with Gaussian-process BO refinement. Across four five-dimensional benchmarks and ten seeds per configuration, every run used exactly 600 objective evaluations. The plain EDA achieved the lowest median final objective on every benchmark. The central research question therefore receives a negative answer for the tested configuration: BO-based local refinement did not improve median evaluation efficiency over allocating the budget to population search.

BO nevertheless compared favorably with the classical hybrid on three of the four marginal medians and won more corresponding-seed comparisons against L-BFGS-B on Sphere, Rastrigin, and Ackley. These results support a limited conclusion about the implemented alternatives, rather than a general superiority claim for BO. The classical optimizer had little budget for numerical gradients, and the two operators differed in their search bounds and reuse of the starting evaluation.

The most useful extension would record the best-so-far objective after every true evaluation, the cost and improvement of each local invocation, and elapsed time. Controlled experiments could then vary local-search frequency and budget, compare equal local domains, reuse nearby observations across BO calls, and adapt the local radius to population spread. A classical baseline with enough evaluations to complete numerical-gradient steps, or with analytic derivatives when available, would clarify whether its current weakness is caused mainly by the ten-call limit. Larger budgets and additional dimensions would test whether the observed ranking persists. Until such evidence is available, the plain Gaussian EDA is the best-supported choice among the three methods for the settings studied here.
